\begin{theorem}[轴外零点扫描剖面的有限秩核化：从边界标量到 \texorpdfstring{$\kappa$}{kappa} 维再现核 Hilbert 空间]\label{thm:xi-basepoint-scan-finite-rank-rkhs}
\leanverified{paper\_xi\_basepoint\_scan\_finite\_rank\_rkhs}
沿用定义 \ref{def:xi-basepoint-scan-profile} 与命题 \ref{prop:xi-basepoint-depth-closed-form} 的记号。
记
\[
a_k:=1+\delta_k,\qquad
p_k:=\gamma_k-\mathrm{i}a_k\in\CC_{-},\qquad
w_k:=4\delta_k\qquad(1\le k\le \kappa).
\]
定义特征映射 \(v:\RR\to\CC^{\kappa}\) 为
\[
v(x)_k:=\frac{\sqrt{w_k}}{x-p_k},
\]
并定义核
\[
K(x,y):=\langle v(x),v(y)\rangle_{\CC^{\kappa}}
=\sum_{k=1}^{\kappa}\frac{w_k}{(x-p_k)(y-\overline{p_k})}
=\sum_{k=1}^{\kappa}\frac{w_k}{(x-\gamma_k+\mathrm{i}a_k)(y-\gamma_k-\mathrm{i}a_k)}.
\]
则 \(K\) 为正定核，且其对角严格等于扫描剖面：
\[
\boxed{\;K(x,x)=H(x)\qquad(x\in\RR).\;}
\]
因此扫描剖面 \(H\) 可被精确核化为一个有限秩核的对角。
并且 \(\mathcal{H}:=\overline{\Span}\{v(x):x\in\RR\}\subseteq\CC^\kappa\) 在自然内积下为 \(\kappa\)-维 Hilbert 空间（等价地：\(K\) 为一维参数生成的 \(\kappa\)-维 RKHS 的再现核）。
\end{theorem}
\begin{proof}
由定义 \(K(x,y)=\langle v(x),v(y)\rangle\)，对任意有限点集 \(\{x_i\}_{i=1}^{n}\subset\RR\) 与系数 \(\{c_i\}_{i=1}^{n}\subset\CC\) 有
\[
\sum_{i,j=1}^{n}c_i\overline{c_j}K(x_i,x_j)
=\left\|\sum_{i=1}^{n}c_i v(x_i)\right\|_{\CC^\kappa}^2\ \ge\ 0,
\]
故 \(K\) 正定。
对角恒等式由 \(|x-p_k|^2=(x-\gamma_k)^2+a_k^2\) 直接给出：
\[
K(x,x)=\sum_{k=1}^{\kappa}\frac{w_k}{|x-p_k|^2}
=\sum_{k=1}^{\kappa}\frac{4\delta_k}{(x-\gamma_k)^2+(1+\delta_k)^2}
=H(x).
\]
最后，为证 \(\dim\mathcal{H}=\kappa\)，取任意 \(\kappa\) 个互异实点 \(x_1,\dots,x_\kappa\)。
矩阵 \(\bigl(\frac{1}{x_i-p_k}\bigr)_{1\le i,k\le\kappa}\) 为 Cauchy 矩阵，其行列式由 Cauchy 恒等式严格非零（因 \(x_i\in\RR\)、\(p_k\in\CC_{-}\) 互不相等），故 \(\{v(x_i)\}_{i=1}^{\kappa}\) 线性无关并张成 \(\CC^\kappa\)。
从而 \(\mathcal{H}=\CC^\kappa\)。
将 \(K\) 视作特征映射 \(v\) 的内积核，即得到相应的 \(\kappa\)-维 RKHS 描述。证毕。
\end{proof}

\begin{proposition}[扫描剖面与 Cauchy 极点骨架的核唯一性]\label{prop:xi-basepoint-scan-kernel-uniqueness}
\leanverified{paper\_xi\_basepoint\_scan\_kernel\_uniqueness}
沿用定理 \ref{thm:xi-basepoint-scan-finite-rank-rkhs} 的记号，并令扫描剖面
\(
H(x)=\sum_{k=1}^{\kappa}\frac{w_k}{|x-p_k|^2}
\)。
设给定矩阵 \(C=(c_{k\ell})\in\CC^{\kappa\times\kappa}\) 并定义核
\[
\widetilde K(x,y):=\sum_{k,\ell=1}^{\kappa}\frac{c_{k\ell}}{(x-p_\ell)(y-\overline{p_k})}\qquad(x,y\in\RR).
\]
若对所有 \(x\in\RR\) 有对角一致性 \(\widetilde K(x,x)=H(x)\)，则必有
\[
C=\mathrm{diag}(w_1,\dots,w_\kappa),
\qquad
\widetilde K(x,y)=\sum_{k=1}^{\kappa}\frac{w_k}{(x-p_k)(y-\overline{p_k})}=K(x,y).
\]
此外，若 \(C\) 为 Hermite 半正定矩阵，则 \(\widetilde K\) 为正定核；当 \(w_k>0\) 时，该核秩为 \(\kappa\)。
\end{proposition}
\begin{proof}
设 \(W:=\mathrm{diag}(w_1,\dots,w_\kappa)\) 并令 \(D:=C-W\)。
由对角一致性，得到对一切 \(x\in\RR\) 的有理恒等式
\[
0=\widetilde K(x,x)-K(x,x)
=\sum_{k,\ell=1}^{\kappa}\frac{D_{k\ell}}{(x-p_\ell)(x-\overline{p_k})}.
\]
右端为关于复变量 \(z\) 的有理函数，因其在实轴上一致为零，故在 \(\CC\) 上恒等为零。
取该有理函数在 \(z=p_j\) 处的留数（\(1\le j\le \kappa\)）得到
\[
0=\operatorname*{Res}_{z=p_j}\sum_{k,\ell=1}^{\kappa}\frac{D_{k\ell}}{(z-p_\ell)(z-\overline{p_k})}
=\sum_{k=1}^{\kappa}\frac{D_{kj}}{p_j-\overline{p_k}}.
\]
令 \(M\in\CC^{\kappa\times\kappa}\) 为 Cauchy 矩阵 \(M_{kj}:=\frac{1}{p_j-\overline{p_k}}\)。
由于 \(\{p_j\}\subset\CC_{-}\) 与 \(\{\overline{p_k}\}\subset\CC_{+}\) 互异，Cauchy 行列式恒等式给出 \(\det M\ne 0\)，故每个列向量 \((D_{kj})_{k}\) 必为零，从而 \(D=0\)，即 \(C=W\)。
\par
最后一段由标准 Gram 表示给出：若 \(C\succeq 0\)，则 \(\widetilde K(x,y)=\langle C^{1/2}u(x),C^{1/2}u(y)\rangle\)（其中 \(u(x)_k=(x-p_k)^{-1}\)）为正定核。
当 \(w_k>0\) 时 \(C=W\) 正定，结合定理 \ref{thm:xi-basepoint-scan-finite-rank-rkhs} 证明中 Cauchy 矩阵非退化性，可得秩为 \(\kappa\)。证毕。
\end{proof}

\begin{proposition}[锚定行列式的 Cauchy--Vandermonde 闭式：满秩锚集下的完全显式分解]\label{prop:xi-basepoint-scan-anchor-det-cauchy-vandermonde}
\leanverified{paper\_xi\_basepoint\_scan\_anchor\_det\_cauchy\_vandermonde}
在定理 \ref{thm:xi-basepoint-scan-finite-rank-rkhs} 的设定下，取互异实锚点集
\(
A=\{x_1,\dots,x_n\}\subset\RR
\)
并令 \(V\in\CC^{n\times\kappa}\) 为
\[
V_{ik}:=v(x_i)_k=\frac{\sqrt{w_k}}{x_i-p_k}.
\]
记其 Gram 矩阵
\[
G_A:=(K(x_i,x_j))_{1\le i,j\le n}=VV^\ast.
\]
则对任意 \(n\le \kappa\) 有严格的 Cauchy--Binet 展开
\[
\boxed{\;
\det G_A
=\sum_{\substack{I\subseteq\{1,\dots,\kappa\}\\ |I|=n}}
\Bigl(\prod_{k\in I}w_k\Bigr)\cdot
\left|\det\!\left(\frac{1}{x_i-p_k}\right)_{\substack{1\le i\le n\\ k\in I}}\right|^2\; }.
\]
特别地，当 \(n=\kappa\)（满秩锚定）时，
\[
\det G_A=\Bigl(\prod_{k=1}^{\kappa}w_k\Bigr)\cdot
\left|\det\!\left(\frac{1}{x_i-p_k}\right)_{1\le i,k\le\kappa}\right|^2
=4^\kappa\Bigl(\prod_{k=1}^{\kappa}\delta_k\Bigr)\,
\frac{|\Delta(X)|^2\,|\Delta(P)|^2}{\prod_{i,k}|x_i-p_k|^2},
\]
其中 \(\Delta(X):=\prod_{1\le i<j\le \kappa}(x_j-x_i)\)，\(\Delta(P):=\prod_{1\le i<j\le \kappa}(p_j-p_i)\)。
\end{proposition}
\begin{proof}
由 \(G_A=VV^\ast\)，对 \(n\le\kappa\) 应用 Cauchy--Binet 得
\[
\det(VV^\ast)
=\sum_{|I|=n}\bigl|\det(V_{(:,I)})\bigr|^2.
\]
并且 \(V_{(:,I)}=\bigl(\frac{1}{x_i-p_k}\bigr)_{k\in I}\cdot \mathrm{diag}(\sqrt{w_k})_{k\in I}\)，故
\(
|\det(V_{(:,I)})|^2
=(\prod_{k\in I}w_k)\cdot\bigl|\det(\frac{1}{x_i-p_k})\bigr|^2
\)，得到第一式。
\par
当 \(n=\kappa\) 时，Cauchy 行列式恒等式给出
\[
\det\!\left(\frac{1}{x_i-p_k}\right)_{1\le i,k\le\kappa}
=\frac{\Delta(X)\Delta(P)}{\prod_{i,k}(x_i-p_k)}.
\]
取模平方并代入 \(w_k=4\delta_k\) 即得结论。证毕。
\end{proof}

\begin{corollary}[满秩锚定行列式的结果式表达：分母等于 \texorpdfstring{$\prod_k|y(p_k)|^2$}{prod|y(pk)|^2}]\label{cor:xi-basepoint-scan-anchor-det-resultant}
\leanverified{paper\_xi\_basepoint\_scan\_anchor\_det\_resultant}
在命题 \ref{prop:xi-basepoint-scan-anchor-det-cauchy-vandermonde} 的满秩情形 \(n=\kappa\) 下，
令锚多项式
\[
y(z):=\prod_{i=1}^{\kappa}(z-x_i)
\qquad(\text{首一，且根两两互异}).
\]
并令极点多项式
\[
\mathcal{A}(z):=\prod_{k=1}^{\kappa}(z-p_k)(z-\overline{p_k})\in\RR[z].
\]
定义结果式
\(
\operatorname{Res}(y,\mathcal{A}):=\prod_{\mathcal{A}(\zeta)=0}y(\zeta)
\)
（按重数计）。
则有严格恒等式
\[
\boxed{\;
\operatorname{Res}(y,\mathcal{A})
=\prod_{k=1}^{\kappa}y(p_k)y(\overline{p_k})
=\prod_{k=1}^{\kappa}|y(p_k)|^{2}\ >\ 0\;}
\]
并且
\[
\boxed{\;
\det G_A
=\Bigl(\prod_{k=1}^{\kappa}w_k\Bigr)\,
\frac{|\Delta(X)|^2\,|\Delta(P)|^2}{\operatorname{Res}(y,\mathcal{A})}\; }.
\]
\end{corollary}
\begin{proof}
由定义，
\(
\operatorname{Res}(y,\mathcal{A})
=\prod_{k=1}^{\kappa}y(p_k)y(\overline{p_k})
\)。
又因 \(y\in\RR[z]\)，有 \(y(\overline{p_k})=\overline{y(p_k)}\)，故该积等于 \(\prod_k|y(p_k)|^2\) 且严格为正（因 \(p_k\notin\RR\) 使 \(y(p_k)\ne 0\)）。
\par
另一方面，由命题 \ref{prop:xi-basepoint-scan-anchor-det-cauchy-vandermonde} 的满秩闭式
\[
\det G_A=\Bigl(\prod_{k=1}^{\kappa}w_k\Bigr)\,
\frac{|\Delta(X)|^2\,|\Delta(P)|^2}{\prod_{i,k}|x_i-p_k|^2}.
\]
注意
\(
\prod_{i,k}|x_i-p_k|^2
=\prod_{k=1}^{\kappa}\prod_{i=1}^{\kappa}|x_i-p_k|^2
=\prod_{k=1}^{\kappa}|y(p_k)|^2
=\operatorname{Res}(y,\mathcal{A})
\)，代回即得。证毕。
\end{proof}

\begin{proposition}[主角分解：锚定互信息的精确几何意义与两点分离度的正弦表示]\label{prop:xi-basepoint-scan-anchor-principal-angles}
\leanverified{paper\_xi\_basepoint\_scan\_anchor\_principal\_angles}
沿用命题 \ref{prop:xi-basepoint-scan-anchor-det-cauchy-vandermonde} 的设定。
对 \(x_i\in A\) 令
\[
u_i:=\frac{v(x_i)}{\|v(x_i)\|},\qquad
\|v(x_i)\|^2=K(x_i,x_i)=H(x_i),
\]
并设相关矩阵 \(R_A:=(\langle u_i,u_j\rangle)_{1\le i,j\le n}\)。
则
\[
\boxed{\;
\det G_A=\Bigl(\prod_{i=1}^{n}H(x_i)\Bigr)\det R_A\; }.
\]
若对 \(\{u_i\}\) 作 Gram--Schmidt 正交化，记 \(\theta_j\in(0,\pi/2]\) 为 \(u_j\) 与 \(\Span\{u_1,\dots,u_{j-1}\}\) 的主角（\(2\le j\le n\)），则
\[
\boxed{\;
\det R_A=\prod_{j=2}^{n}\sin^2\theta_j,\qquad
\log\det G_A=\sum_{i=1}^{n}\log H(x_i)+2\sum_{j=2}^{n}\log(\sin\theta_j)\; }.
\]
特别地，当 \(n=2\) 时，
\[
\det R_A=1-|\langle u_1,u_2\rangle|^2=\sin^2\theta_2.
\]
\end{proposition}
\begin{proof}
由 \(v(x_i)=\|v(x_i)\|u_i=\sqrt{H(x_i)}\,u_i\) 得 \(G_A=DR_AD\)，其中 \(D=\mathrm{diag}(\sqrt{H(x_1)},\dots,\sqrt{H(x_n)})\)。
取行列式即得第一式。
\par
对 \(\{u_i\}\) 作 Gram--Schmidt，记 \(d_j\) 为 \(u_j\) 到 \(\Span\{u_1,\dots,u_{j-1}\}\) 的正交距离，则 \(\det R_A=\prod_{j=1}^{n}d_j^2\)（Gram 行列式等于正交化后对角元素平方乘积），且 \(d_1=1\)、\(d_j=\sin\theta_j\)（\(j\ge 2\)）。
代入即得所示分解；两点情形为直接计算。证毕。
\end{proof}

\begin{corollary}[近极大锚定的角度集中：行列式亏损控制“非正交事件”的比例]\label{cor:xi-basepoint-scan-anchor-angle-concentration}
沿用命题 \ref{prop:xi-basepoint-scan-anchor-principal-angles} 的记号。
若
\[
\det G_A\ \ge\ e^{-\varepsilon}\prod_{i=1}^{n}H(x_i)\qquad(\varepsilon\ge 0),
\]
等价于 \(\det R_A\ge e^{-\varepsilon}\)，则对任意 \(\eta\in(0,1)\)，至少有 \((1-\eta)(n-1)\) 个指标 \(j\in\{2,\dots,n\}\) 满足
\[
\boxed{\;
\sin^2\theta_j\ \ge\ \exp\!\Bigl(-\frac{\varepsilon}{\eta(n-1)}\Bigr)\; }.
\]
\end{corollary}
\begin{proof}
由命题 \ref{prop:xi-basepoint-scan-anchor-principal-angles}，有 \(\prod_{j=2}^{n}\sin^2\theta_j=\det R_A\ge e^{-\varepsilon}\)，取对数得
\[
\sum_{j=2}^{n}-\log(\sin^2\theta_j)\le \varepsilon.
\]
若有超过 \(\eta(n-1)\) 个 \(j\) 满足 \(-\log(\sin^2\theta_j)>\varepsilon/(\eta(n-1))\)，则上式左端严格大于 \(\varepsilon\)，矛盾。
故除至多 \(\eta(n-1)\) 个异常指标外，其余皆满足所示下界。证毕。
\end{proof}

\begin{proposition}[\texorpdfstring{$\mathrm{PSL}(2,\RR)$}{PSL(2,R)} 共变性：扫描核与锚定行列式的严格 Möbius 变换律]\label{prop:xi-basepoint-scan-psl2r-covariance}
在定理 \ref{thm:xi-basepoint-scan-finite-rank-rkhs} 的设定下，取 Möbius 变换
\[
m(z)=\frac{az+b}{cz+d}\qquad(a,b,c,d\in\RR,\ ad-bc>0),
\]
其将上半平面保向映到自身。
令 \(x_i':=m(x_i)\) 与 \(p_k':=m(p_k)\)。
定义变换后的权重
\[
w_k':=w_k\cdot\frac{(ad-bc)^2}{|cp_k+d|^2}.
\]
以 \(\{(p_k',w_k')\}\) 生成变换后特征映射 \(v'\)、核 \(K'\) 与 Gram 矩阵 \(G_{A'}'\)。
则对所有 \(x,y\in\RR\) 有严格共变关系
\[
\boxed{\;
K'(m(x),m(y))=(cx+d)(cy+d)\,K(x,y)\; }.
\]
从而对 \(A=\{x_1,\dots,x_n\}\) 有
\[
\boxed{\;
\det G_{A'}'=\Bigl(\prod_{i=1}^{n}(cx_i+d)^2\Bigr)\det G_A\; }.
\]
并且归一化相关矩阵 \(R_A\)（命题 \ref{prop:xi-basepoint-scan-anchor-principal-angles}）在该变换下不变，从而由 \(\det R_A\) 与主角序列 \(\{\theta_j\}\) 所定义的“纯角度信息”是 \(\mathrm{PSL}(2,\RR)\)-不变量。
\end{proposition}
\begin{proof}
恒等式
\[
\frac{1}{m(x)-m(p)}=\frac{(cx+d)(cp+d)}{ad-bc}\cdot\frac{1}{x-p}
\]
对一切 \(x\neq p\) 成立。
取 \(p=p_k\) 并乘以 \(\sqrt{w_k'}\) 得
\[
\frac{\sqrt{w_k'}}{m(x)-m(p_k)}
=(cx+d)\cdot \frac{cp_k+d}{|cp_k+d|}\cdot \frac{\sqrt{w_k}}{x-p_k}.
\]
因此存在对角酉矩阵 \(U=\mathrm{diag}\bigl((cp_k+d)/|cp_k+d|\bigr)\) 使对一切 \(x\in\RR\) 有
\(
v'(m(x))=(cx+d)\,U\,v(x)
\)。
取内积并用 \(U\) 的酉性即得核的共变式；Gram 行列式缩放由 \(G_{A'}'=D G_A D\)（\(D=\mathrm{diag}(cx_i+d)\)）立得。
最后，相关矩阵在左右公共缩放下不变，故主角与 \(\det R_A\) 不变。证毕。
\end{proof}

\begin{proposition}[满秩锚定的驻点方程：实轴电荷--复极点平衡的显式系统]\label{prop:xi-basepoint-scan-full-rank-anchor-critical-equations}
在命题 \ref{prop:xi-basepoint-scan-anchor-det-cauchy-vandermonde} 的满秩情形 \(n=\kappa\) 下，
将 \(\det G_A\) 视为互异实变量 \((x_1,\dots,x_\kappa)\) 的函数。
若 \(A\) 为开域内的临界点（对每个 \(x_i\) 的偏导为零），则每个 \(i=1,\dots,\kappa\) 满足平衡方程
\[
\boxed{\;
\sum_{\substack{j=1\\ j\ne i}}^{\kappa}\frac{1}{x_i-x_j}
=\sum_{k=1}^{\kappa}\Re\frac{1}{x_i-p_k}
=\sum_{k=1}^{\kappa}\frac{x_i-\gamma_k}{(x_i-\gamma_k)^2+a_k^2}\; }.
\]
\end{proposition}
\begin{proof}
由命题 \ref{prop:xi-basepoint-scan-anchor-det-cauchy-vandermonde} 的闭式，
\[
\log\det G_A=\mathrm{const}+2\log|\Delta(X)|-2\sum_{i,k}\log|x_i-p_k|,
\]
其中常数项与 \(\{p_k,w_k\}\) 无关。
对 \(x_i\) 求导得
\(
\partial_{x_i}\log|\Delta(X)|=\sum_{j\ne i}\frac{1}{x_i-x_j}
\)
且
\(
\partial_{x_i}\log|x_i-p_k|=\Re\frac{1}{x_i-p_k}
\)。
令 \(\partial_{x_i}\log\det G_A=0\) 即得第一条等式；最后一式为实部展开。证毕。
\end{proof}

\leanverified{paper\_xi\_basepoint\_scan\_heine\_stieltjes\_ode}
\begin{proposition}[锚定临界构型的 Heine--Stieltjes 型微分方程：Van Vleck 多项式的存在唯一性]\label{prop:xi-basepoint-scan-heine-stieltjes-ode}
在命题 \ref{prop:xi-basepoint-scan-full-rank-anchor-critical-equations} 的满秩情形 \(n=\kappa\) 下，
设锚点 \(X=\{x_1,\dots,x_\kappa\}\subset\RR\) 两两互异，并令锚多项式
\(
y(z):=\prod_{j=1}^{\kappa}(z-x_j)
\)
为首一多项式。
令极点多项式
\[
\mathcal{A}(z):=\prod_{k=1}^{\kappa}(z-p_k)(z-\overline{p_k})\in\RR[z].
\]
则以下两者等价：
\begin{enumerate}
\item 锚点构型 \(X\) 为 \(\det G_A\) 的临界构型（等价地满足命题 \ref{prop:xi-basepoint-scan-full-rank-anchor-critical-equations} 中的平衡方程）。
\item 存在多项式 \(B\in\RR[z]\)（\(\deg B\le 2\kappa-2\)）使恒等式
\[
\boxed{\;
\mathcal{A}(z)\,y''(z)-\mathcal{A}'(z)\,y'(z)=B(z)\,y(z)\;}
\]
在 \(\CC\) 上成立。
\end{enumerate}
并且当该恒等式成立时，多项式 \(B\) 由 \(y\) 唯一确定（称为对应的 Van Vleck 多项式）。
\end{proposition}
\begin{proof}
记
\(
N(z):=\mathcal{A}(z)y''(z)-\mathcal{A}'(z)y'(z)
\)。
\par
\emph{(1)\(\Rightarrow\)(2).}
对任意根 \(x_i\)（故 \(y(x_i)=0\)、且 \(y'(x_i)\ne 0\)）有恒等式
\[
\frac{y''(x_i)}{y'(x_i)}=2\sum_{j\ne i}\frac{1}{x_i-x_j}.
\]
另一方面，
\[
\frac{\mathcal{A}'(z)}{\mathcal{A}(z)}
=\sum_{k=1}^{\kappa}\left(\frac{1}{z-p_k}+\frac{1}{z-\overline{p_k}}\right),
\]
从而当 \(z=x_i\in\RR\) 时
\(
\frac{\mathcal{A}'(x_i)}{\mathcal{A}(x_i)}
=2\sum_{k=1}^{\kappa}\Re\frac{1}{x_i-p_k}
\)。
将其与临界平衡方程
\(
\sum_{j\ne i}\frac{1}{x_i-x_j}=\sum_{k}\Re\frac{1}{x_i-p_k}
\)
合并，得到
\(
y''(x_i)/y'(x_i)=\mathcal{A}'(x_i)/\mathcal{A}(x_i)
\)，等价于
\[
N(x_i)=\mathcal{A}(x_i)y''(x_i)-\mathcal{A}'(x_i)y'(x_i)=0.
\]
由于 \(\{x_i\}\) 两两互异且 \(y\) 为首一的 \(\kappa\) 次多项式，故 \(y\mid N\)。
令 \(B:=N/y\)，则 \(B\in\RR[z]\)，且
\(
\deg B\le \deg N-\deg y\le (2\kappa+\kappa-2)-\kappa=2\kappa-2
\)。
\par
\emph{(2)\(\Rightarrow\)(1).}
设 \(\mathcal{A}y''-\mathcal{A}'y'=By\) 成立。
对任意根 \(x_i\) 代入得
\(
\mathcal{A}(x_i)y''(x_i)-\mathcal{A}'(x_i)y'(x_i)=0
\)。
又因
\(
\mathcal{A}(x_i)=\prod_{k}|x_i-p_k|^2>0
\)，可除得
\(
y''(x_i)/y'(x_i)=\mathcal{A}'(x_i)/\mathcal{A}(x_i)
\)。
将 \(\mathcal{A}'/\mathcal{A}\) 的部分分式展开代入，并用
\(
y''(x_i)/y'(x_i)=2\sum_{j\ne i}\frac{1}{x_i-x_j}
\)，即得临界平衡方程
\(
\sum_{j\ne i}\frac{1}{x_i-x_j}=\sum_k\Re\frac{1}{x_i-p_k}
\)。
由命题 \ref{prop:xi-basepoint-scan-full-rank-anchor-critical-equations}，这等价于 \(X\) 为 \(\det G_A\) 的临界构型。
\par
最后，若 \(\mathcal{A}y''-\mathcal{A}'y'=B_1y=B_2y\)，则 \((B_1-B_2)y\equiv 0\) 强制 \(B_1=B_2\)。
故 \(B\) 唯一。证毕。
\end{proof}

\begin{proposition}[Van Vleck 多项式的首项刚性与次首项闭式]\label{prop:xi-basepoint-scan-van-vleck-leading-coeffs}
在命题 \ref{prop:xi-basepoint-scan-heine-stieltjes-ode} 的设定下，设 \(y\) 首一且 \(\deg y=\kappa\)，并令
\[
B(z)=b_0 z^{2\kappa-2}+b_1 z^{2\kappa-3}+O(z^{2\kappa-4}),
\quad
\mathcal{A}(z)=z^{2\kappa}+\mathcal{A}_1 z^{2\kappa-1}+O(z^{2\kappa-2}),
\quad
y(z)=z^{\kappa}+y_1 z^{\kappa-1}+O(z^{\kappa-2}).
\]
则必有
\[
\boxed{\;b_0=-\kappa(\kappa+1),\qquad b_1=-\kappa^{2}\mathcal{A}_1+2y_1.\;}
\]
其中
\(
\mathcal{A}_1=-\sum_{k=1}^{\kappa}(p_k+\overline{p_k})=-2\sum_{k=1}^{\kappa}\gamma_k
\)，且
\(
y_1=-\sum_{j=1}^{\kappa}x_j
\)。
\end{proposition}
\begin{proof}
将恒等式 \(\mathcal{A}y''-\mathcal{A}'y'=By\) 在无穷远处作首项比较。
由
\[
y'(z)=\kappa z^{\kappa-1}+(\kappa-1)y_1 z^{\kappa-2}+O(z^{\kappa-3}),
\qquad
y''(z)=\kappa(\kappa-1)z^{\kappa-2}+(\kappa-1)(\kappa-2)y_1 z^{\kappa-3}+O(z^{\kappa-4}),
\]
以及
\[
\mathcal{A}'(z)=2\kappa z^{2\kappa-1}+(2\kappa-1)\mathcal{A}_1 z^{2\kappa-2}+O(z^{2\kappa-3}),
\]
得到
\[
\mathcal{A}y''-\mathcal{A}'y'
=\bigl(\kappa(\kappa-1)-2\kappa^2\bigr)z^{3\kappa-2}
\bigl(-\kappa^2\mathcal{A}_1-(\kappa-1)(\kappa+2)y_1\bigr)z^{3\kappa-3}
+O(z^{3\kappa-4}).
\]
另一方面，
\[
By=b_0 z^{3\kappa-2}+(b_1+b_0y_1)z^{3\kappa-3}+O(z^{3\kappa-4}).
\]
逐项比较即得 \(b_0=-\kappa(\kappa+1)\) 与 \(b_1=-\kappa^{2}\mathcal{A}_1+2y_1\)。证毕。
\end{proof}

\begin{proposition}[Van Vleck 比值的零/一阶矩残差律：缺陷电流的严格线性约束]\label{prop:xi-basepoint-scan-van-vleck-residue-moment01}
在命题 \ref{prop:xi-basepoint-scan-heine-stieltjes-ode} 的设定下，令
\(
f(z):=\frac{B(z)}{\mathcal{A}(z)}
\)。
则 \(f\) 为仅在 \(\{p_k,\overline{p_k}\}\) 处有一阶极点的有理函数，且由命题 \ref{prop:xi-basepoint-scan-van-vleck-leading-coeffs} 有
\(
f(z)=-\kappa(\kappa+1)z^{-2}+O(z^{-3})
\)。
记
\[
r_k:=\operatorname*{Res}_{z=p_k} f(z)=\frac{B(p_k)}{\mathcal{A}'(p_k)}.
\]
则必有
\[
\boxed{\;
\sum_{k=1}^{\kappa}\Re r_k=0,\qquad
\sum_{k=1}^{\kappa}\Re(p_k r_k)=-\frac{\kappa(\kappa+1)}{2}\; }.
\]
若 \(\mathcal{A}y''-\mathcal{A}'y'=By\) 成立，则进一步在每个 \(p_k\) 处有
\[
r_k=-\frac{y'(p_k)}{y(p_k)},
\]
从而等价地可写为
\[
\boxed{\;
\sum_{k=1}^{\kappa}\Re\frac{y'(p_k)}{y(p_k)}=0,\qquad
\sum_{k=1}^{\kappa}\Re\!\left(p_k\frac{y'(p_k)}{y(p_k)}\right)=\frac{\kappa(\kappa+1)}{2}\; }.
\]
\end{proposition}
\begin{proof}
由于 \(\deg B=\deg\mathcal{A}-2\)，\(f\) 在无穷远处无多项式部分。
又因 \(B,\mathcal{A}\in\RR[z]\)，故 \(f\) 的极点成共轭对，且其部分分式展开可写为
\[
f(z)=\sum_{k=1}^{\kappa}\frac{r_k}{z-p_k}+\sum_{k=1}^{\kappa}\frac{\overline{r_k}}{z-\overline{p_k}}.
\]
展开于无穷远得到
\[
f(z)=\frac{\sum_k(r_k+\overline{r_k})}{z}
\;+\;\frac{\sum_k(p_k r_k+\overline{p_k}\,\overline{r_k})}{z^2}
\;+\;O(z^{-3}).
\]
而由命题 \ref{prop:xi-basepoint-scan-van-vleck-leading-coeffs}，
\(
f(z)=-\kappa(\kappa+1)z^{-2}+O(z^{-3})
\)，其 \(z^{-1}\) 系数为零，故 \(\sum_k\Re r_k=0\)。
其 \(z^{-2}\) 系数为 \(-\kappa(\kappa+1)\)，故
\(
2\sum_k\Re(p_k r_k)=-\kappa(\kappa+1)
\)。
\par
最后，若 \(\mathcal{A}y''-\mathcal{A}'y'=By\) 成立，则在 \(z=p_k\) 处因 \(\mathcal{A}(p_k)=0\) 得
\[
-\mathcal{A}'(p_k)y'(p_k)=B(p_k)y(p_k),
\]
从而 \(r_k=B(p_k)/\mathcal{A}'(p_k)=-y'(p_k)/y(p_k)\)。证毕。
\end{proof}

\begin{proposition}[二阶矩残差律与“重心项”显式：极点平方加权的闭式恒等式]\label{prop:xi-basepoint-scan-van-vleck-residue-moment2}
沿用命题 \ref{prop:xi-basepoint-scan-van-vleck-residue-moment01} 的记号。
则有严格恒等式
\[
\boxed{\;
\sum_{k=1}^{\kappa}\Re(p_k^{2} r_k)=\frac{1}{2}\bigl(b_1-b_0\mathcal{A}_1\bigr)\; }.
\]
在 \(r_k=-y'(p_k)/y(p_k)\) 的情形下，结合命题 \ref{prop:xi-basepoint-scan-van-vleck-leading-coeffs} 的
\(
b_0=-\kappa(\kappa+1)
\)、
\(
b_1=-\kappa^{2}\mathcal{A}_1+2y_1
\)，并用 \(\mathcal{A}_1=-2\sum_{k}\gamma_k\)、\(y_1=-\sum_{j}x_j\) 化简，得到
\[
\boxed{\;
\sum_{k=1}^{\kappa}\Re\!\left(p_k^{2}\frac{y'(p_k)}{y(p_k)}\right)
=\kappa\sum_{k=1}^{\kappa}\gamma_k+\sum_{j=1}^{\kappa}x_j\; }.
\]
\end{proposition}
\begin{proof}
由命题 \ref{prop:xi-basepoint-scan-van-vleck-residue-moment01} 的部分分式展开，展开于无穷远得其 \(z^{-3}\) 系数为
\[
\sum_{k=1}^{\kappa}(p_k^2 r_k+\overline{p_k}^2\,\overline{r_k})
=2\sum_{k=1}^{\kappa}\Re(p_k^2 r_k).
\]
另一方面，
\[
\frac{B(z)}{\mathcal{A}(z)}
=b_0 z^{-2}+(b_1-b_0\mathcal{A}_1)z^{-3}+O(z^{-4}),
\]
比较 \(z^{-3}\) 系数即得第一式。
\par
若 \(r_k=-y'(p_k)/y(p_k)\)，则将其与命题 \ref{prop:xi-basepoint-scan-van-vleck-leading-coeffs} 的 \(b_0,b_1\) 代入第一式，并用 \(\mathcal{A}_1=-2\sum\gamma_k\)、\(y_1=-\sum x_j\) 化简，即得第二式。证毕。
\end{proof}

\begin{proposition}[满秩锚定行列式的全局极大值存在性与临界点必然性]\label{prop:xi-basepoint-scan-anchor-det-global-maximum}
在命题 \ref{prop:xi-basepoint-scan-anchor-det-cauchy-vandermonde} 的满秩情形 \(n=\kappa\) 下，
视 \(\det G_A\) 为实配置空间
\[
\Omega:=\{(x_1,\dots,x_\kappa)\in\RR^\kappa:\ x_i\neq x_j\}
\]
上的连续函数（对置换对称）。
则 \(\det G_A\) 存在全局极大值；任取极大点 \(X^\star\in\Omega\)，其必为临界构型，从而满足命题 \ref{prop:xi-basepoint-scan-full-rank-anchor-critical-equations} 中的平衡方程。
进而由命题 \ref{prop:xi-basepoint-scan-heine-stieltjes-ode}，存在对应的首一锚多项式 \(y^\star\) 与唯一的 Van Vleck 多项式 \(B^\star\) 使
\[
\mathcal{A}(z)\,(y^\star)''(z)-\mathcal{A}'(z)\,(y^\star)'(z)=B^\star(z)\,y^\star(z)
\]
成立。
\leanverified{paper\_xi\_basepoint\_scan\_anchor\_det\_global\_maximum}
\end{proposition}
\begin{proof}
由定理 \ref{thm:xi-basepoint-scan-finite-rank-rkhs}，对任意互异锚点 \(X\in\Omega\)，Gram 矩阵 \(G_A\) 正定，故 \(\det G_A>0\)。
令 \(M:=\sup_{X\in\Omega}\det G_A\in(0,\infty)\)，并取极大序列 \(X^{(m)}\in\Omega\) 使 \(\det G_{A^{(m)}}\to M\)。
\par
首先排除无穷远逃逸。
由 Hadamard 不等式与 \(G_{ii}=H(x_i)\)（定理 \ref{thm:xi-basepoint-scan-finite-rank-rkhs}），有
\[
\det G_A\le \prod_{i=1}^{\kappa}H(x_i).
\]
而 \(H(x)=\sum_k w_k/|x-p_k|^2=O(x^{-2})\)（\(|x|\to\infty\)），故若 \(\max_i|x_i^{(m)}|\to\infty\)，则右端趋于 \(0\)，与 \(\det G_{A^{(m)}}\to M>0\) 矛盾。
因此 \(\{X^{(m)}\}\) 有界，从而存在收敛子列（仍记为 \(X^{(m)}\)）使 \(X^{(m)}\to X^\star\in\RR^\kappa\)。
\par
其次排除碰撞边界。
由命题 \ref{prop:xi-basepoint-scan-anchor-det-cauchy-vandermonde} 的满秩闭式，
\(
\det G_A\propto |\Delta(X)|^2
\)；
当发生碰撞 \(x_i\to x_j\) 时 \(\Delta(X)\to 0\)，故 \(\det G_A\to 0\)。
因此极大序列的极限点 \(X^\star\) 不可能落在碰撞边界上，必满足 \(X^\star\in\Omega\)。
由连续性得 \(\det G_{A^\star}=M\)，即全局极大值达到。
\par
由于 \(\Omega\) 为开集，极大点 \(X^\star\) 为内点，故 \(\nabla \det G_A(X^\star)=0\)，即为临界构型。
由命题 \ref{prop:xi-basepoint-scan-full-rank-anchor-critical-equations} 得平衡方程成立，最后由命题 \ref{prop:xi-basepoint-scan-heine-stieltjes-ode} 得到相应的微分方程与 \(B^\star\) 的存在唯一性。证毕。
\end{proof}

\begin{proposition}[锚定临界构型的解析刚性：极点参数的实解析依赖与局部孤立性]\label{prop:xi-basepoint-scan-critical-configuration-analytic-rigidity}
在命题 \ref{prop:xi-basepoint-scan-full-rank-anchor-critical-equations} 的满秩情形下，设极点多重集 \(\{p_k\}\subset\CC_{-}\) 互异。
考虑映射 \(\Phi=(\Phi_1,\dots,\Phi_\kappa)\)：
\[
\Phi_i(x_1,\dots,x_\kappa;\,p_1,\dots,p_\kappa)
:=\sum_{j\ne i}\frac{1}{x_i-x_j}-\sum_{k=1}^{\kappa}\Re\frac{1}{x_i-p_k},
\qquad i=1,\dots,\kappa.
\]
若存在临界构型 \(X^\star=(x_1^\star,\dots,x_\kappa^\star)\) 满足 \(\Phi(X^\star;p)=0\)，且该点的 Jacobian
\(
J:=\bigl(\partial \Phi_i/\partial x_j\bigr)_{1\le i,j\le \kappa}
\)
可逆，则在 \(p\) 的充分小邻域内存在唯一实解析分支 \(X(p)\) 使 \(\Phi(X(p);p)=0\) 且 \(X(p^\star)=X^\star\)。
因此在非退化条件下，临界锚定构型对极点参数呈现实解析刚性，并在该邻域内局部孤立。
\end{proposition}
\begin{proof}
\(\Phi\) 在互异点域上对 \((x,p)\) 为实解析函数。
可逆 Jacobian 假设使隐函数定理可用，故存在唯一实解析解分支 \(X(p)\) 并给出局部孤立性。证毕。
\end{proof}

\paragraph{Van Vleck 系数矩锁定、临界 Jacobian 的谱化与实平移协变}
在锚定临界构型的 Heine--Stieltjes 型微分方程、Van Vleck 留数矩约束以及极点参数的局部解析刚性已经建立之后，还可继续抽取出四条彼此闭合的后果。它们分别把 Van Vleck 的前两项系数压缩为极点--残差实矩、把非退化条件具体化为一个实对称谱判据、把临界分支上的一阶灵敏度写成无须求解 \(\partial X/\partial p\) 的 envelope 核，以及把共同实平移下的 \(b_1\) 漂移分解为显式仿射项与严格不变量。

\leanverified{paper\_xi\_basepoint\_scan\_van\_vleck\_coeff\_residue\_bilinear\_lock}
\begin{theorem}[Van Vleck 前两项系数与极点--残差实矩的双线性锁定]\label{thm:xi-basepoint-scan-van-vleck-coeff-residue-bilinear-lock}
沿用命题 \ref{prop:xi-basepoint-scan-van-vleck-leading-coeffs} 与
\ref{prop:xi-basepoint-scan-van-vleck-residue-moment2} 的记号。则成立
\[
b_0=2\sum_{k=1}^{\kappa}\Re(p_k r_k),
\]
以及
\[
b_1
=
2\sum_{k=1}^{\kappa}\Re(p_k^2 r_k)-\kappa(\kappa+1)\mathcal A_1
=
2\sum_{k=1}^{\kappa}\Re(p_k^2 r_k)
+2\kappa(\kappa+1)\sum_{k=1}^{\kappa}\gamma_k.
\]
因此，Van Vleck 多项式的前两项系数完全由极点--残差的一、二阶实矩决定；其中 \(b_0\) 仅依赖于阶数 \(\kappa\)，而 \(b_1\) 对极点实部质心只表现为严格线性漂移。
\end{theorem}
\begin{proof}
由命题 \ref{prop:xi-basepoint-scan-van-vleck-residue-moment01}，
\[
2\sum_{k=1}^{\kappa}\Re(p_k r_k)=-\kappa(\kappa+1),
\]
而命题 \ref{prop:xi-basepoint-scan-van-vleck-leading-coeffs} 给出
\(
b_0=-\kappa(\kappa+1)
\)，故得第一式。

再由命题 \ref{prop:xi-basepoint-scan-van-vleck-residue-moment2}，
\[
2\sum_{k=1}^{\kappa}\Re(p_k^2 r_k)=b_1-b_0\mathcal A_1.
\]
将 \(b_0=-\kappa(\kappa+1)\) 代入，得到
\[
b_1
=
2\sum_{k=1}^{\kappa}\Re(p_k^2 r_k)-\kappa(\kappa+1)\mathcal A_1.
\]
最后用
\(
\mathcal A_1=-2\sum_{k=1}^{\kappa}\gamma_k
\)
即得第二种写法。证毕。
\end{proof}

\leanverified{paper\_xi\_basepoint\_scan\_critical\_jacobian\_calogero\_symmetric}
\begin{theorem}[非退化临界锚定构型的 Jacobian 为实对称 Calogero 型矩阵]\label{thm:xi-basepoint-scan-critical-jacobian-calogero-symmetric}
沿用命题 \ref{prop:xi-basepoint-scan-critical-configuration-analytic-rigidity} 的记号。对 \(\Phi=(\Phi_1,\dots,\Phi_\kappa)\) 的定义域内任意互异实配置 \(X=(x_1,\dots,x_\kappa)\)，其 Jacobian
\[
J:=\left(\frac{\partial\Phi_i}{\partial x_j}\right)_{1\le i,j\le \kappa}
\]
满足显式公式
\[
J_{ij}=
\begin{cases}
\displaystyle \frac{1}{(x_i-x_j)^2}, & i\neq j,\\[3mm]
\displaystyle -\sum_{\ell\neq i}\frac{1}{(x_i-x_\ell)^2}
+\sum_{k=1}^{\kappa}\Re\frac{1}{(x_i-p_k)^2}, & i=j.
\end{cases}
\]
特别地，
\[
J=J^{\mathsf T}\in M_\kappa(\RR).
\]
因此，在命题 \ref{prop:xi-basepoint-scan-critical-configuration-analytic-rigidity} 的非退化口径下，“Jacobian 可逆”等价于实对称矩阵 \(J\) 不含零特征值。
\end{theorem}
\begin{proof}
对 \(i\neq j\)，只有和式
\(
\sum_{m\neq i}(x_i-x_m)^{-1}
\)
中的 \(m=j\) 项依赖 \(x_j\)，并且
\[
\frac{\partial}{\partial x_j}\frac{1}{x_i-x_j}
=\frac{1}{(x_i-x_j)^2},
\]
故
\[
J_{ij}=\frac{1}{(x_i-x_j)^2}.
\]

对对角元，
\[
\frac{\partial}{\partial x_i}\sum_{j\neq i}\frac{1}{x_i-x_j}
=-\sum_{j\neq i}\frac{1}{(x_i-x_j)^2}.
\]
另一方面，由
\[
\Phi_i=
\sum_{j\neq i}\frac{1}{x_i-x_j}
-\sum_{k=1}^{\kappa}\Re\frac{1}{x_i-p_k}
\]
可知
\[
\frac{\partial}{\partial x_i}\left(-\Re\frac{1}{x_i-p_k}\right)
=\Re\frac{1}{(x_i-p_k)^2}.
\]
两式相加即得 \(J_{ii}\) 的公式。

由于对 \(i\neq j\) 有 \(J_{ij}=J_{ji}\)，而对角元显然为实数，故 \(J\) 为实对称矩阵。于是其可逆性等价于 \(0\notin\sigma(J)\)。证毕。
\end{proof}

\begin{theorem}[临界分支归一化极值行列式的 envelope 公式]\label{thm:xi-basepoint-scan-normalized-envelope-derivative-formula}
沿用命题 \ref{prop:xi-basepoint-scan-critical-configuration-analytic-rigidity} 的设定。设 \(p^\star\) 处存在非退化临界构型，并记 \(X(p)=(x_1(p),\dots,x_\kappa(p))\) 为该命题给出的唯一实解析分支。令
\[
A(p):=\{x_1(p),\dots,x_\kappa(p)\},
\qquad
P:=(p_1,\dots,p_\kappa),
\]
并定义局部极值值函数
\[
M(p):=\det G_{A(p)}.
\]
再定义去除极点前因子的归一化极值函数
\[
\widetilde M(p)
:=
\frac{M(p)}{\left(\prod_{k=1}^{\kappa}w_k\right)|\Delta(P)|^2}
=
\frac{|\Delta(X(p))|^2}{\prod_{i,k}|x_i(p)-p_k|^2}.
\]
则 \(\widetilde M(p)\) 为实解析函数，并且对每个 \(\ell=1,\dots,\kappa\) 都有
\[
\frac{\partial}{\partial \gamma_\ell}\log \widetilde M(p)
=
2\sum_{i=1}^{\kappa}
\frac{x_i(p)-\gamma_\ell}{(x_i(p)-\gamma_\ell)^2+a_\ell^2},
\]
\[
\frac{\partial}{\partial a_\ell}\log \widetilde M(p)
=
-2\sum_{i=1}^{\kappa}
\frac{a_\ell}{(x_i(p)-\gamma_\ell)^2+a_\ell^2}.
\]
换言之，沿临界分支的一阶灵敏度不含 \(\partial X/\partial p\) 项，而完全由边界势核的显式和式给出。
等价地，
\[
\log M(p)=\log\!\Bigl(\left(\prod_{k=1}^{\kappa}w_k\right)|\Delta(P)|^2\Bigr)+\log\widetilde M(p),
\]
故 \(\log M(p)\) 的导数只比上式多出由极点前因子 \(\left(\prod_k w_k\right)|\Delta(P)|^2\) 贡献的显式项。
\end{theorem}
\begin{proof}
由命题 \ref{prop:xi-basepoint-scan-critical-configuration-analytic-rigidity}，分支 \(X(p)\) 为实解析，故 \(\widetilde M(p)\) 亦为实解析函数。

由命题 \ref{prop:xi-basepoint-scan-anchor-det-cauchy-vandermonde} 的满秩闭式，
\[
\log \widetilde M(p)
=
2\log|\Delta(X(p))|
-2\sum_{i,k}\log|x_i(p)-p_k|.
\]
对任一参数 \(\eta\in\{\gamma_\ell,a_\ell\}\) 作链式求导，得到
\[
\frac{\partial}{\partial \eta}\log \widetilde M(p)
=
\sum_{j=1}^{\kappa}
\frac{\partial}{\partial x_j}\log \widetilde M(p)\cdot \frac{\partial x_j}{\partial \eta}
+\left(\frac{\partial}{\partial \eta}\log \widetilde M(p)\right)_{\!X\ \mathrm{fixed}}.
\]
由于归一化前因子仅依赖极点而不依赖 \(X\)，故
\[
\frac{\partial}{\partial x_j}\log \widetilde M(p)
=
\frac{\partial}{\partial x_j}\log \det G_{A(p)}.
\]
而沿临界分支，命题 \ref{prop:xi-basepoint-scan-full-rank-anchor-critical-equations} 给出
\(
\nabla_X \log \det G_A=0
\)，
故所有含 \(\partial X/\partial \eta\) 的项全部消失。

因此只需计算显式极点导数。对 \(\gamma_\ell\)，有
\[
\frac{\partial}{\partial \gamma_\ell}\log|x_i-p_\ell|
=-\frac{x_i-\gamma_\ell}{(x_i-\gamma_\ell)^2+a_\ell^2},
\]
故
\[
\frac{\partial}{\partial \gamma_\ell}\log \widetilde M(p)
=
-2\sum_{i=1}^{\kappa}\frac{\partial}{\partial \gamma_\ell}\log|x_i-p_\ell|
=
2\sum_{i=1}^{\kappa}
\frac{x_i(p)-\gamma_\ell}{(x_i(p)-\gamma_\ell)^2+a_\ell^2}.
\]
同理，
\[
\frac{\partial}{\partial a_\ell}\log|x_i-p_\ell|
=\frac{a_\ell}{(x_i-\gamma_\ell)^2+a_\ell^2},
\]
故
\[
\frac{\partial}{\partial a_\ell}\log \widetilde M(p)
=
-2\sum_{i=1}^{\kappa}
\frac{a_\ell}{(x_i(p)-\gamma_\ell)^2+a_\ell^2}.
\]
证毕。
\end{proof}

\begin{theorem}[共同实平移下的临界协变与 Van Vleck 次首项的仿射转移律]\label{thm:xi-basepoint-scan-real-translation-b1-affine-transfer}
沿用命题 \ref{prop:xi-basepoint-scan-full-rank-anchor-critical-equations}、\ref{prop:xi-basepoint-scan-critical-configuration-analytic-rigidity} 与命题 \ref{prop:xi-basepoint-scan-van-vleck-leading-coeffs} 的记号。对任意 \(t\in\RR\)，记
\[
p+t:=(p_1+t,\dots,p_\kappa+t).
\]
则成立：
\begin{enumerate}
  \item 若 \(X=(x_1,\dots,x_\kappa)\) 是极点组 \(p\) 的临界构型，则
  \[
  X+t:=(x_1+t,\dots,x_\kappa+t)
  \]
  是极点组 \(p+t\) 的临界构型。
  \item 若 \(X(p)\) 为命题 \ref{prop:xi-basepoint-scan-critical-configuration-analytic-rigidity} 给出的局部唯一实解析分支，则对所有使 \(p+t\) 仍落在其定义域内的 \(t\)，有
  \[
  X(p+t)=X(p)+t.
  \]
  \item 设 \(B_p(z)=b_0(p)z^{2\kappa-2}+b_1(p)z^{2\kappa-3}+O(z^{2\kappa-4})\) 为对应的 Van Vleck 多项式，则
  \[
  b_0(p+t)=b_0(p),
  \qquad
  b_1(p+t)=b_1(p)+2\kappa(\kappa^2-1)t.
  \]
\end{enumerate}
因此
\[
\mathcal I(p)
:=
b_1(p)+(\kappa^2-1)\mathcal A_1(p)
=
b_1(p)-2(\kappa^2-1)\sum_{k=1}^{\kappa}\gamma_k
\]
是共同实平移下的严格不变量。
\end{theorem}
\begin{proof}
临界方程
\[
\sum_{j\ne i}\frac{1}{x_i-x_j}
=
\sum_{k=1}^{\kappa}\Re\frac{1}{x_i-p_k}
\]
只依赖于差值 \(x_i-x_j\) 与 \(x_i-p_k\)。因此同时把全部 \(x_i\) 与全部 \(p_k\) 加上同一实数 \(t\) 时，方程保持不变，从而第 \(1\) 条成立。

对第 \(2\) 条，固定 \(t\) 并把平移后的参数记为 \(q\)。则
\[
q\longmapsto X(q-t)+t
\]
是参数 \(q\) 的一条解析解分支。由命题 \ref{prop:xi-basepoint-scan-critical-configuration-analytic-rigidity} 的局部唯一性，在两条分支共同定义的区域上只能有
\[
X(q)=X(q-t)+t.
\]
令 \(q=p+t\) 即得
\(
X(p+t)=X(p)+t
\)。

对第 \(3\) 条，由命题 \ref{prop:xi-basepoint-scan-van-vleck-leading-coeffs}，
\[
b_0=-\kappa(\kappa+1),
\qquad
b_1=-\kappa^2\mathcal A_1+2y_1.
\]
共同实平移下，
\[
\mathcal A_1(p+t)=\mathcal A_1(p)-2\kappa t,
\qquad
y_1(p+t)=y_1(p)-\kappa t.
\]
代回得
\[
b_1(p+t)
=
-\kappa^2\bigl(\mathcal A_1(p)-2\kappa t\bigr)+2\bigl(y_1(p)-\kappa t\bigr)
=
b_1(p)+2\kappa(\kappa^2-1)t,
\]
而 \(b_0\) 本身与极点位置无关，故保持不变。

最后，
\[
\mathcal I(p+t)
=
b_1(p+t)+(\kappa^2-1)\mathcal A_1(p+t)
\]
中由 \(b_1\) 产生的增量 \(2\kappa(\kappa^2-1)t\) 与由 \(\mathcal A_1\) 产生的增量 \(-2\kappa(\kappa^2-1)t\) 恰好抵消，故 \(\mathcal I(p+t)=\mathcal I(p)\)。证毕。
\end{proof}

\endinput


% (User input window)
\leanverified{paper\_xi\_basepoint\_scan\_kernel\_completeness\_full\_rank}

\paragraph{有限秩锚定核的完备关系与 Gram 逆闭式：由离散锚点完全重构连续核}
沿用定理 \ref{thm:xi-basepoint-scan-finite-rank-rkhs} 的记号：极点 \(\{p_k\}_{k=1}^{\kappa}\subset\CC_{-}\)、权 \(\{w_k\}_{k=1}^{\kappa}\subset\RR_{>0}\)、特征映射 \(v\) 与核 \(K\)。

\begin{proposition}[满秩锚集的核完备关系：由离散锚点完全重构连续核]\label{prop:xi-basepoint-scan-kernel-completeness-full-rank}
对任意 \(x\in\RR\) 定义
\[
v(x)_k:=\frac{\sqrt{w_k}}{x-p_k}\in\CC,\qquad
K(x,y):=\langle v(x),v(y)\rangle_{\CC^\kappa}
:=\sum_{k=1}^{\kappa}\frac{w_k}{(x-p_k)(y-\overline{p_k})}\qquad(x,y\in\RR).
\]
取互异实锚点 \(A=\{a_1,\dots,a_\kappa\}\subset\RR\)，并记 Gram 矩阵
\[
G_A:=(K(a_i,a_j))_{1\le i,j\le \kappa}.
\]
则 \(G_A\) 可逆，且对一切 \(x,y\in\RR\) 成立完备恒等式
\[
\boxed{\;
K(x,y)=K(x,A)\,G_A^{-1}\,K(A,y)\; },
\]
其中
\(
K(x,A)=(K(x,a_1),\dots,K(x,a_\kappa))
\)、
\(
K(A,y)=(K(a_1,y),\dots,K(a_\kappa,y))^{\mathsf T}
\)。
\end{proposition}
\begin{proof}
令矩阵 \(V\in\CC^{\kappa\times\kappa}\) 为
\(
V_{ik}:=v(a_i)_k=\sqrt{w_k}/(a_i-p_k)
\)。
则 \(G_A=VV^\ast\)。
由定理 \ref{thm:xi-basepoint-scan-finite-rank-rkhs} 证明中 Cauchy 非退化性可知 \(V\) 可逆，
故
\[
I_\kappa=V^\ast G_A^{-1}V.
\]
两边左乘 \(v(x)^\ast\)、右乘 \(v(y)\) 得
\[
K(x,y)
:=\langle v(x),v(y)\rangle
:=\langle v(x),V^\ast G_A^{-1}Vv(y)\rangle
:=K(x,A)\,G_A^{-1}\,K(A,y),
\]
即所示完备关系。证毕。
\end{proof}

\begin{proposition}[满秩 Gram 逆的全闭式：由 Cauchy 逆显式给出 \texorpdfstring{$G_A^{-1}$}{G\_A^{-1}}]\label{prop:xi-basepoint-scan-gram-inverse-closed-form}
\leanverified{paper\_xi\_basepoint\_scan\_gram\_inverse\_closed\_form}
在命题 \ref{prop:xi-basepoint-scan-kernel-completeness-full-rank} 的满秩设定下，令
\[
\Pi(z):=\prod_{k=1}^{\kappa}(z-p_k),\qquad
y(z):=\prod_{j=1}^{\kappa}(z-a_j).
\]
记 \(V_{ik}=\sqrt{w_k}/(a_i-p_k)\)，则 \(V=C\,\mathrm{diag}(\sqrt w)\)，其中 \(C_{ik}=(a_i-p_k)^{-1}\)。
由结论 \ref{con:xi-terminal-cauchy-lagrange-rational-interpolation-cauchy-inverse} 的 Cauchy 逆闭式可得
\[
(V^{-1})_{k i}
:=\frac{1}{\sqrt{w_k}}\,(C^{-1})_{k i}
:=\frac{1}{\sqrt{w_k}}\cdot
\frac{y(p_k)\,\Pi(a_i)}{(p_k-a_i)\,\Pi'(p_k)\,y'(a_i)}.
\]
从而
\[
G_A^{-1}=(V^{-1})^\ast V^{-1},
\]
并且其元素具有显式表达：对 \(1\le i,j\le \kappa\),
\[
\boxed{\;
(G_A^{-1})_{ij}
:=
\frac{\overline{\Pi(a_i)}\,\Pi(a_j)}{y'(a_i)\,y'(a_j)}
\sum_{k=1}^{\kappa}
\frac{|y(p_k)|^2}{w_k\,|\Pi'(p_k)|^2}\cdot
\frac{1}{(a_i-\overline{p_k})(a_j-p_k)}\; }.
\]
\end{proposition}
\begin{proof}
由 \(G_A=VV^\ast\) 且 \(V\) 可逆，得 \(G_A^{-1}=(V^{-1})^\ast V^{-1}\)。
而 \(V=C\,\mathrm{diag}(\sqrt w)\) 给出 \(V^{-1}=\mathrm{diag}(w^{-1/2})\,C^{-1}\)，代入结论 \ref{con:xi-terminal-cauchy-lagrange-rational-interpolation-cauchy-inverse} 的 \(C^{-1}\) 闭式即得 \((V^{-1})_{ki}\)。
将
\(
(G_A^{-1})_{ij}=\sum_{k}\overline{(V^{-1})_{ki}}(V^{-1})_{kj}
\)
逐项代入并整理得到所示表达式。证毕。
\end{proof}

% (User input window)
%
% Full-rank anchoring: weight gauge invariance and explicit Λ-formulas.

\begin{proposition}[满秩锚定的权规范不变性：插值系数的权无关性与 \texorpdfstring{$D$}{D}-最优锚构型的权不变性]\label{prop:xi-basepoint-scan-full-rank-weight-gauge-invariance}
\leanverified{paper\_xi\_basepoint\_scan\_full\_rank\_weight\_gauge\_invariance}
在命题 \ref{prop:xi-basepoint-scan-kernel-completeness-full-rank} 的满秩设定下，定义插值系数向量
\[
\Lambda_A(x):=G_A^{-1}k_A(x)\in\CC^\kappa,
\qquad
k_A(x):=(K(a_1,x),\dots,K(a_\kappa,x))^{\mathsf T}.
\]
则成立：
\begin{enumerate}
  \item \(\Lambda_A(x)\) 与权重 \(\{w_k\}\) 无关；更精确地，令 \(C(A;P)=(\frac{1}{a_i-p_k})_{i,k}\) 与 \(u(x)=(\frac{1}{x-p_k})_{k}\)，则
  \[
  \boxed{\;\Lambda_A(x)=\bigl(C(A;P)^\ast\bigr)^{-1}u(x)\;}
  \]
  其中右端仅依赖 \((A,\{p_k\})\)。
  \item \(\det G_A\) 的权因子分离为
  \[
  \boxed{\;\det G_A=\left(\prod_{k=1}^{\kappa}w_k\right)\cdot \left|\det C(A;P)\right|^2.\;}
  \]
  因此在极点集合 \(\{p_k\}\) 固定而权重任意变化时，最大化 \(\det G_A\) 的满秩锚构型集合（\(D\)-最优设计）不随权重改变。
\end{enumerate}
\end{proposition}
\begin{proof}
令 \(W^{1/2}:=\mathrm{diag}(\sqrt{w_1},\dots,\sqrt{w_\kappa})\)，并记
\[
V=(v(a_i)_k)_{i,k}=\left(\frac{\sqrt{w_k}}{a_i-p_k}\right)_{i,k}=C(A;P)\,W^{1/2}.
\]
则 \(G_A=VV^\ast\)，并由 Cauchy 非退化性知 \(V\) 可逆。
从而
\[
\det G_A=\det(VV^\ast)=|\det V|^2
=|\det C(A;P)|^2\,|\det W^{1/2}|^2
=\left(\prod_{k=1}^{\kappa}w_k\right)\cdot|\det C(A;P)|^2,
\]
得到第二条，且 \(D\)-最优锚构型的权不变性立即推出。
\par
另一方面，由 \(k_A(x)=Vv(x)\)（参见命题 \ref{prop:xi-basepoint-scan-gram-determinant-schur-increment} 证明中的恒等式）与 \(G_A=VV^\ast\)，有
\[
\Lambda_A(x)=G_A^{-1}k_A(x)=(VV^\ast)^{-1}Vv(x)=(V^\ast)^{-1}v(x).
\]
又 \(V^\ast=W^{1/2}C(A;P)^\ast\)，而 \(v(x)=W^{1/2}u(x)\)。
故 \(\Lambda_A(x)=(C(A;P)^\ast)^{-1}u(x)\)，其中不含 \(\{w_k\}\)。证毕。
\end{proof}

\begin{proposition}[满秩插值系数的 Cauchy--Lagrange 闭式]\label{prop:xi-basepoint-scan-lambda-cauchy-lagrange-closed-form}
在命题 \ref{prop:xi-basepoint-scan-full-rank-weight-gauge-invariance} 的设定下，令
\[
y(z):=\prod_{j=1}^{\kappa}(z-a_j),\qquad \Pi(z):=\prod_{k=1}^{\kappa}(z-p_k),
\]
并记 \(\Lambda_A(x)=(\Lambda_{A,i}(x))_{i=1}^{\kappa}\)。
则对每个 \(1\le i\le \kappa\) 有显式表达式
\[
\boxed{\;
\Lambda_{A,i}(x)
=
-\frac{\Pi(a_i)}{y'(a_i)}
\sum_{k=1}^{\kappa}
\frac{y(p_k)}{(x-p_k)(a_i-p_k)\,\Pi'(p_k)}\; }.
\]
特别地，该闭式与权重 \(\{w_k\}\) 无关。
\end{proposition}
\begin{proof}
由命题 \ref{prop:xi-basepoint-scan-full-rank-weight-gauge-invariance}，\(\Lambda_A(x)\) 等价于求解线性系统
\[
\sum_{i=1}^{\kappa}\frac{\Lambda_{A,i}(x)}{a_i-p_k}=\frac{1}{x-p_k}\qquad(1\le k\le \kappa),
\]
其系数矩阵为 Cauchy 矩阵 \(C(A;P)\)。
令
\(
P(z):=\sum_{i=1}^{\kappa}\Lambda_{A,i}(x)\,\frac{y(z)}{z-a_i}
\)
则 \(\deg P\le \kappa-1\) 且满足 \(\Lambda_{A,i}(x)=P(a_i)/y'(a_i)\)。
将上式在 \(z=p_k\) 处取值得
\[
P(p_k)=\frac{y(p_k)}{p_k-x}\qquad(1\le k\le \kappa).
\]
因此 \(P\) 为结点 \(\{p_k\}\) 上对数据 \(y(p_k)/(p_k-x)\) 的 Lagrange 插值多项式：
\[
P(z)=\sum_{k=1}^{\kappa}\frac{y(p_k)}{p_k-x}\cdot \frac{\Pi(z)}{(z-p_k)\Pi'(p_k)}.
\]
令 \(z=a_i\) 得
\(
P(a_i)=\Pi(a_i)\sum_{k}\frac{y(p_k)}{(p_k-x)(a_i-p_k)\Pi'(p_k)}
\)。
再用 \((p_k-x)^{-1}=-(x-p_k)^{-1}\) 并代回 \(\Lambda_{A,i}(x)=P(a_i)/y'(a_i)\) 即得所示闭式。证毕。
\end{proof}

\endinput


\begin{proposition}[行列式增量与正交缺口：Schur 补的几何意义]\label{prop:xi-basepoint-scan-gram-determinant-schur-increment}
\leanverified{paper\_xi\_basepoint\_scan\_gram\_determinant\_schur\_increment}
对任意锚集 \(A=\{a_1,\dots,a_n\}\subset\RR\)（\(n<\kappa\)）假设 \(G_A:=(K(a_i,a_j))_{i,j}\) 可逆，并定义
\[
k_A(x):=(K(x,a_1),\dots,K(x,a_n))^{\mathsf T},\qquad
H(x):=K(x,x).
\]
则对任意 \(x\in\RR\setminus A\) 有严格恒等式
\[
\boxed{\;
\det G_{A\cup\{x\}}=\det G_A\cdot\Bigl(H(x)-k_A(x)^\ast G_A^{-1}k_A(x)\Bigr)\; } .
\]
并且括号内量非负，且满足正交缺口表示
\[
H(x)-k_A(x)^\ast G_A^{-1}k_A(x)
:=\bigl\|v(x)-\mathrm{proj}_{\Span\{v(a_i)\}}v(x)\bigr\|_{\CC^\kappa}^{2}.
\]
\end{proposition}
\begin{proof}
写成分块矩阵
\[
G_{A\cup\{x\}}=
\begin{pmatrix}
G_A & k_A(x)\\
k_A(x)^\ast & H(x)
\end{pmatrix}.
\]
对分块行列式用 Schur 补公式即得第一式。
第二式由 \(G_A=V_AV_A^\ast\)（\(V_A\) 为由 \(v(a_i)\) 组成的 \(n\times\kappa\) 矩阵）与最小二乘投影恒等式推出：
\(
k_A(x)=V_A v(x)
\)，从而 Schur 补等于 \(v(x)\) 在 \(\Span\{v(a_i)\}\) 的正交缺口平方范数。证毕。
\end{proof}

\begin{corollary}[Gram 行列式的杠杆分解：有限秩核下的乘法因子化与全局上界]\label{cor:xi-basepoint-scan-gram-determinant-leverage-product}
\leanverified{paper\_xi\_basepoint\_scan\_gram\_determinant\_leverage\_product}
对任意有序锚列 \((a_1,\dots,a_n)\)（互异，且每步 \(G_{A_m}\) 可逆），记递增锚集 \(A_m:=\{a_1,\dots,a_m\}\)（\(A_0=\varnothing\)），并定义杠杆缺口
\[
\ell_{A_{m-1}}(a_m):=
H(a_m)-k_{A_{m-1}}(a_m)^\ast G_{A_{m-1}}^{-1}k_{A_{m-1}}(a_m)\ \ge 0.
\]
则成立严格乘法分解
\[
\boxed{\;
\det G_{A_n}=\prod_{m=1}^{n}\ell_{A_{m-1}}(a_m)\; }.
\]
从而对任意 \(n\) 点锚集 \(A\)（忽略顺序）有上界
\[
\det G_A\le \prod_{m=1}^{n}\Bigl(\sup_{x\in\RR}\ell_{A_{m-1}}(x)\Bigr),
\]
且当且仅当存在某个排列使得每一步 \(a_m\) 都实现相应的上确界时取等。
\end{corollary}
\begin{proof}
由命题 \ref{prop:xi-basepoint-scan-gram-determinant-schur-increment} 对 \(m=1,\dots,n\) 逐步应用得
\[
\det G_{A_m}=\det G_{A_{m-1}}\cdot \ell_{A_{m-1}}(a_m),
\]
连乘即得分解。
上界由每因子 \(\ell_{A_{m-1}}(a_m)\le \sup_x\ell_{A_{m-1}}(x)\) 逐项相乘得到；
取等条件由每步取等的必要充分性给出。证毕。
\end{proof}

\begin{theorem}[余维一锚定的杠杆缺口闭式：权重与形状的完全分离]\label{thm:xi-basepoint-scan-leverage-gap-codim1-closed-form}
\leanverified{paper\_xi\_basepoint\_scan\_leverage\_gap\_codim1\_closed\_form}
沿用命题 \ref{prop:xi-basepoint-scan-gram-determinant-schur-increment} 的记号，并假设 \(|A|=\kappa-1\) 且 \(G_A\) 可逆。
令
\[
\Pi(z):=\prod_{k=1}^{\kappa}(z-p_k),
\qquad
y_A(z):=\prod_{a\in A}(z-a).
\]
则对任意 \(x\in\RR\setminus A\) 有严格恒等式
\[
\boxed{\;
\ell_A(x)
:=
\frac{\bigl|\frac{y_A(x)}{\Pi(x)}\bigr|^2}{
\displaystyle\sum_{k=1}^{\kappa}\frac{|y_A(p_k)|^2}{w_k\,|\Pi'(p_k)|^2}
}\; }.
\]
特别地，\(\ell_A(x)\) 作为 \(x\) 的函数，其形状仅由 \((A,\{p_k\})\) 决定；权重 \(\{w_k\}\) 仅通过整体比例常数
\(
\Bigl(\sum_{k=1}^{\kappa}\frac{|y_A(p_k)|^2}{w_k|\Pi'(p_k)|^2}\Bigr)^{-1}
\)
进入。
\end{theorem}
\begin{proof}
由命题 \ref{prop:xi-basepoint-scan-gram-determinant-schur-increment}，
\(
\ell_A(x)=\| \mathrm{proj}_{\Span\{v(a):a\in A\}^\perp}v(x)\|^2
\)。
当 \(|A|=\kappa-1\) 时，\(\Span\{v(a):a\in A\}^\perp\) 为一维。
令
\[
d_k:=\frac{y_A(p_k)}{\Pi'(p_k)},
\qquad
c_k:=\frac{\overline{d_k}}{\sqrt{w_k}}
\qquad(1\le k\le \kappa),
\]
并记 \(c=(c_k)_{k=1}^{\kappa}\in\CC^\kappa\)。
对任意 \(a\in A\)，有
\[
\langle v(a),c\rangle
:=\sum_{k=1}^{\kappa}\frac{\sqrt{w_k}}{a-p_k}\cdot\overline{c_k}
:=\sum_{k=1}^{\kappa}\frac{d_k}{a-p_k}.
\]
另一方面，由部分分式恒等式
\[
\frac{y_A(z)}{\Pi(z)}
:=\sum_{k=1}^{\kappa}\frac{y_A(p_k)}{\Pi'(p_k)}\cdot\frac{1}{z-p_k}
:=\sum_{k=1}^{\kappa}\frac{d_k}{z-p_k},
\]
代入 \(z=a\in A\) 得
\(
\sum_k\frac{d_k}{a-p_k}=\frac{y_A(a)}{\Pi(a)}=0
\)，故 \(c\perp v(a)\) 对所有 \(a\in A\) 成立，从而
\(
\Span\{v(a):a\in A\}^\perp=\Span\{c\}
\)。
于是由一维投影公式
\[
\ell_A(x)=\frac{|\langle v(x),c\rangle|^2}{\|c\|^2}.
\]
再由同一部分分式恒等式，
\[
\langle v(x),c\rangle
:=\sum_{k=1}^{\kappa}\frac{\sqrt{w_k}}{x-p_k}\cdot\overline{c_k}
:=\sum_{k=1}^{\kappa}\frac{d_k}{x-p_k}
:=\frac{y_A(x)}{\Pi(x)}.
\]
并且
\(
\|c\|^2=\sum_{k=1}^{\kappa}|c_k|^2=\sum_{k=1}^{\kappa}\frac{|d_k|^2}{w_k}
:=\sum_{k=1}^{\kappa}\frac{|y_A(p_k)|^2}{w_k|\Pi'(p_k)|^2}
\)。
代回即得结论。证毕。
\end{proof}

\begin{corollary}[余维一锚定的临界点方程：与权重无关的平衡式]\label{cor:xi-basepoint-scan-leverage-gap-codim1-critical-points}
在定理 \ref{thm:xi-basepoint-scan-leverage-gap-codim1-closed-form} 的假设下，对任意 \(x\in\RR\setminus A\) 有
\[
\boxed{\;
\frac{\dd}{\dd x}\log \ell_A(x)
:=
2\sum_{a\in A}\frac{1}{x-a}
-2\,\Re\sum_{k=1}^{\kappa}\frac{1}{x-p_k}\; }.
\]
因此 \(\ell_A\) 的临界点全体满足
\[
\boxed{\;
\sum_{a\in A}\frac{1}{x-a}
:=
\Re\sum_{k=1}^{\kappa}\frac{1}{x-p_k}\; },
\]
并且该方程与权重 \(\{w_k\}\) 无关。
\end{corollary}
\begin{proof}
由定理 \ref{thm:xi-basepoint-scan-leverage-gap-codim1-closed-form}，
\(
\ell_A(x)=C(A,P,w)\,\bigl|\frac{y_A(x)}{\Pi(x)}\bigr|^2
\)
且 \(C(A,P,w)\) 与 \(x\) 无关。
对实变量 \(x\) 与可微复值函数 \(f\) 有恒等式
\(
\frac{\dd}{\dd x}\log|f(x)|^2=2\Re\bigl(f'(x)/f(x)\bigr)
\)。
又
\[
\frac{y_A'(x)}{y_A(x)}=\sum_{a\in A}\frac{1}{x-a},
\qquad
\frac{\Pi'(x)}{\Pi(x)}=\sum_{k=1}^{\kappa}\frac{1}{x-p_k}.
\]
合并即得。证毕。
\end{proof}

\begin{theorem}[一般余维下的杠杆缺口矩阵闭式：由极点离散矩完全控制]\label{thm:xi-basepoint-scan-leverage-gap-general-codim-matrix}
沿用命题 \ref{prop:xi-basepoint-scan-gram-determinant-schur-increment} 的记号，并设 \(|A|=n<\kappa\) 且 \(G_A\) 可逆。
记 \(m:=\kappa-n\ge 1\)，仍令 \(\Pi(z):=\prod_{k=1}^{\kappa}(z-p_k)\)、\(y_A(z):=\prod_{a\in A}(z-a)\)。
对 \(r=0,1,\dots,m-1\) 定义向量 \(c^{(r)}\in\CC^\kappa\)：
\[
c^{(r)}_k:=\frac{\overline{p_k^{\,r}y_A(p_k)/\Pi'(p_k)}}{\sqrt{w_k}}
\qquad(1\le k\le\kappa).
\]
令
\[
C_A:=[c^{(0)}\ \cdots\ c^{(m-1)}]\in\CC^{\kappa\times m},
\qquad
M_A:=C_A^\ast C_A\in\CC^{m\times m},
\qquad
r_A(x):=C_A^\ast v(x)\in\CC^{m}.
\]
则 \(M_A\) 正定，且对任意 \(x\in\RR\setminus A\) 有
\[
\boxed{\;
\ell_A(x)=r_A(x)^\ast M_A^{-1}r_A(x)\; }.
\]
并且 \(r_A(x)\) 具有完全显式形式
\[
\boxed{\;
r_A(x)
:=\left(\frac{y_A(x)}{\Pi(x)},\ \frac{x\,y_A(x)}{\Pi(x)},\ \dots,\ \frac{x^{m-1}y_A(x)}{\Pi(x)}\right)^{\mathsf T}\; }.
\]
等价地，\(\ell_A(x)\) 为有限维正定二次型在向量 \(r_A(x)\) 上的取值，其中离散矩阵元素为
\[
(M_A)_{rs}
:=\sum_{k=1}^{\kappa}\frac{p_k^{\,r}\overline{p_k^{\,s}}\ |y_A(p_k)|^2}{w_k\,|\Pi'(p_k)|^2}
\qquad(0\le r,s\le m-1).
\]
\end{theorem}
\begin{proof}
对每个多项式 \(q\)（\(\deg q\le m-1\)），由部分分式恒等式
\[
\frac{q(z)\,y_A(z)}{\Pi(z)}
:=\sum_{k=1}^{\kappa}\frac{q(p_k)\,y_A(p_k)}{\Pi'(p_k)}\cdot\frac{1}{z-p_k}
\]
并代入 \(z=a\in A\) 得到右端为零，因而向量
\(
d(q)_k:=q(p_k)\,y_A(p_k)/\Pi'(p_k)
\)
满足 \(\sum_k d(q)_k/(a-p_k)=0\) 对所有 \(a\in A\) 成立。
这意味着 \(\ker(V_A)\)（其中 \(V_A\) 为以 \(v(a)^\ast\) 为行的矩阵）包含由 \(\{c^{(r)}\}_{r=0}^{m-1}\) 张成的 \(m\)-维子空间；
而由定理 \ref{thm:xi-basepoint-scan-finite-rank-rkhs} 给出的 \(\dim\Span\{v(x):x\in\RR\}=\kappa\) 与 \(|A|=n\) 可得 \(\dim\ker(V_A)=m\)，故
\(\ker(V_A)=\Span\{c^{(0)},\dots,c^{(m-1)}\}\)。
于是正交投影算子满足
\(
\mathrm{proj}_{\ker(V_A)}=C_A(M_A)^{-1}C_A^\ast
\)，从而
\[
\ell_A(x)
:=\|\mathrm{proj}_{\ker(V_A)}v(x)\|^2
:=v(x)^\ast C_A(M_A)^{-1}C_A^\ast v(x)
:=r_A(x)^\ast M_A^{-1}r_A(x).
\]
最后，对 \(r=0,\dots,m-1\),
\[
(r_A(x))_r
:=\langle c^{(r)},v(x)\rangle
:=\sum_{k=1}^{\kappa}\frac{p_k^{\,r}y_A(p_k)}{\Pi'(p_k)}\cdot\frac{1}{x-p_k}
:=\frac{x^r y_A(x)}{\Pi(x)},
\]
其中最后一步取 \(q(z)=z^r\) 即得。证毕。
\end{proof}

\begin{theorem}[一般余维杠杆缺口的普适 \texorpdfstring{$x^{-2}$}{x^-2} 尾律]\label{thm:xi-basepoint-scan-leverage-gap-universal-xminus2-tail}
在定理 \ref{thm:xi-basepoint-scan-leverage-gap-general-codim-matrix} 的设定下，记 \(e_m\in\CC^m\) 为最后一个标准基向量，并定义
\[
c_A:=e_m^\ast M_A^{-1}e_m>0.
\]
则当 \(|x|\to\infty\) 时，杠杆缺口满足普适渐近
\[
\boxed{\;
\ell_A(x)=\frac{c_A}{x^2}+O(x^{-3}).\;}
\]
特别地，该衰减指数恒等于 \(2\)，与余维 \(m=\kappa-|A|\)、锚点个数 \(|A|\) 以及极点总数 \(\kappa\) 无关。
\end{theorem}
\begin{proof}
由定理 \ref{thm:xi-basepoint-scan-leverage-gap-general-codim-matrix}，
\[
\ell_A(x)=r_A(x)^\ast M_A^{-1}r_A(x),
\qquad
r_A(x)=\frac{y_A(x)}{\Pi(x)}(1,x,\dots,x^{m-1})^{\mathsf T}.
\]
又因 \(\deg y_A=\kappa-m\)、\(\deg\Pi=\kappa\)，故
\[
\frac{y_A(x)}{\Pi(x)}=x^{-m}\bigl(1+O(x^{-1})\bigr)
\qquad(|x|\to\infty).
\]
于是
\[
r_A(x)
=
\bigl(x^{-m},x^{1-m},\dots,x^{-1}\bigr)^{\mathsf T}\bigl(1+O(x^{-1})\bigr)
=
x^{-1}e_m+O(x^{-2}).
\]
代入二次型得到
\[
\ell_A(x)
=
\bigl(x^{-1}e_m+O(x^{-2})\bigr)^\ast
M_A^{-1}
\bigl(x^{-1}e_m+O(x^{-2})\bigr)
=
x^{-2}e_m^\ast M_A^{-1}e_m+O(x^{-3}),
\]
即所求渐近式。
由于 \(M_A\) 正定，故其逆矩阵亦正定，进而 \(c_A=e_m^\ast M_A^{-1}e_m>0\)。
证毕。
\end{proof}

\begin{corollary}[严格非退化性：任取不超过 \texorpdfstring{$\kappa$}{kappa} 个互异锚点均给出可逆 Gram]\label{cor:xi-basepoint-scan-gram-nondegeneracy-up-to-kappa}
在定理 \ref{thm:xi-basepoint-scan-finite-rank-rkhs} 的设定下，若 \(A\subset\RR\) 为互异锚点族且 \(|A|\le\kappa\)，则 \(G_A\) 自动可逆。
并且当 \(|A|<\kappa\) 时，对任意 \(x\in\RR\setminus A\) 有 \(\ell_A(x)>0\)，从而
\[
\det G_{A\cup\{x\}}=\det G_A\cdot \ell_A(x)\ >\ 0.
\]
\end{corollary}
\begin{proof}
取 \(|A|=n\le\kappa\)。从 \(\RR\setminus A\) 中任选 \(\kappa-n\) 个点补齐为 \(\kappa\) 个互异实点 \(A'\supseteq A\)。
由定理 \ref{thm:xi-basepoint-scan-finite-rank-rkhs} 的证明，\(\{v(a):a\in A'\}\) 线性无关，从而其任意子集（特别地 \(\{v(a):a\in A\}\)）亦线性无关。
因此 \(G_A\) 作为 Gram 矩阵为正定并可逆。
\par
若 \(|A|<\kappa\) 且 \(x\notin A\)，同理可把 \(A\cup\{x\}\) 补齐为 \(\kappa\) 点集，从而 \(\{v(a):a\in A\cup\{x\}\}\) 线性无关。
于是 \(v(x)\notin\Span\{v(a):a\in A\}\)，由命题 \ref{prop:xi-basepoint-scan-gram-determinant-schur-increment} 的正交缺口表示得 \(\ell_A(x)>0\)。
最后的行列式正性由 Schur 补增量公式给出。证毕。
\end{proof}

\begin{proposition}[\texorpdfstring{$\mathrm{PSL}(2,\RR)$}{PSL(2,R)} 共变性下的杠杆缺口变换律与比值不变量]\label{prop:xi-basepoint-scan-leverage-gap-psl2r-invariant-ratio}
在命题 \ref{prop:xi-basepoint-scan-psl2r-covariance} 的设定下，取任意锚集 \(A\subset\RR\)（\(|A|<\kappa\)）与 \(x\in\RR\setminus A\)。
令 \(A':=m(A)\)、\(x':=m(x)\)，并以共变重标定后的数据定义 \(\ell_{A'}'(x')\)、\(H'(x')\)。
则有严格变换律
\[
\boxed{\;
\ell_{A'}'(x')=(cx+d)^2\,\ell_A(x),
\qquad
H'(x')=(cx+d)^2\,H(x)\; }.
\]
从而归一化杠杆比值
\[
\boxed{\;
\frac{\ell_A(x)}{H(x)}
\quad\text{在}\ \mathrm{PSL}(2,\RR)\ \text{作用下不变。}\;}
\]
\end{proposition}
\begin{proof}
由命题 \ref{prop:xi-basepoint-scan-psl2r-covariance}，对任意 \(u,v\in\RR\) 有
\(
K'(m(u),m(v))=(cu+d)(cv+d)K(u,v)
\)。
取 \(u=v=x\) 得 \(H'(x')=(cx+d)^2H(x)\)。
同时对锚点 \(A=\{a_i\}_{i=1}^{n}\) 写 \(D=\mathrm{diag}(ca_i+d)\)，则 \(G_{A'}'=DG_A D\)，且 \(k_{A'}'(x')=(cx+d)\,D\,k_A(x)\)。
代入 \(\ell_{A'}'(x')=H'(x')-k_{A'}'(x')^\ast (G_{A'}')^{-1}k_{A'}'(x')\) 即得 \(\ell_{A'}'(x')=(cx+d)^2\ell_A(x)\)。
比值不变量立得。证毕。
\end{proof}

% (User input window)
%
% Residual kernel and polynomial frames (general codimension).

\begin{proposition}[余维一锚定的条件余核：秩一闭式与 \texorpdfstring{$\ell_A=\mathcal{R}_A(\cdot,\cdot)$}{ellA=residual}]\label{prop:xi-basepoint-scan-residual-kernel-codim1-rankone}
\leanverified{paper\_xi\_basepoint\_scan\_residual\_kernel\_codim1\_rankone}
在定理 \ref{thm:xi-basepoint-scan-leverage-gap-codim1-closed-form} 的设定下（\(|A|=\kappa-1\) 且 \(G_A\) 可逆），定义条件余核
\[
\mathcal{R}_A(x,y):=K(x,y)-K(x,A)\,G_A^{-1}\,K(A,y)\qquad(x,y\in\RR).
\]
则 \(\mathcal{R}_A\) 为秩一正定核，并满足闭式
\[
\boxed{\;
\mathcal{R}_A(x,y)
=
\frac{\displaystyle \frac{y_A(x)\,\overline{y_A(y)}}{\Pi(x)\,\overline{\Pi(y)}}}{
\displaystyle \sum_{k=1}^{\kappa}\frac{|y_A(p_k)|^2}{w_k\,|\Pi'(p_k)|^2}
}\; }.
\]
特别地，对角即为杠杆缺口：
\[
\boxed{\;\ell_A(x)=\mathcal{R}_A(x,x)\; }.
\]
\end{proposition}
\begin{proof}
由命题 \ref{prop:xi-basepoint-scan-gram-determinant-schur-increment} 的投影解释，
\[
\mathcal{R}_A(x,y)=\left\langle \mathrm{proj}_{\Span\{v(a):a\in A\}^\perp}v(x),\ \mathrm{proj}_{\Span\{v(a):a\in A\}^\perp}v(y)\right\rangle,
\]
其中 \(\Span\{v(a):a\in A\}^\perp\) 为一维（因 \(|A|=\kappa-1\)）。
定理 \ref{thm:xi-basepoint-scan-leverage-gap-codim1-closed-form} 证明中构造的向量
\(
c_k=\overline{y_A(p_k)/\Pi'(p_k)}/\sqrt{w_k}
\)
张成该正交补，并且 \(\langle v(x),c\rangle=y_A(x)/\Pi(x)\)、\(\|c\|^2=\sum_k |y_A(p_k)|^2/(w_k|\Pi'(p_k)|^2)\)。
由一维投影公式
\(
\mathrm{proj}_{\Span\{c\}}v(x)=\langle v(x),c\rangle\,c/\|c\|^2
\)
代入即得所示闭式与秩一性；取 \(x=y\) 得 \(\ell_A(x)=\mathcal{R}_A(x,x)\)。证毕。
\end{proof}

\begin{theorem}[一般余维下的余核分解：残差多项式框架与正交归一性]\label{thm:xi-basepoint-scan-residual-kernel-polynomial-frame}
\leanverified{paper\_xi\_basepoint\_scan\_residual\_kernel\_polynomial\_frame}
沿用命题 \ref{prop:xi-basepoint-scan-gram-determinant-schur-increment} 的记号，设 \(|A|=n\le\kappa\) 且 \(G_A\) 可逆，并记余维数 \(r:=\kappa-n\ge 1\)。
定义条件余核
\[
\mathcal{R}_A(x,y):=K(x,y)-K(x,A)\,G_A^{-1}\,K(A,y)\qquad(x,y\in\RR).
\]
则存在 \(r\) 个次数 \(\le \kappa-1\) 的多项式 \(y_{A,1},\dots,y_{A,r}\) 使：
\begin{enumerate}
  \item 余核具有规范分解
  \[
  \boxed{\;
  \mathcal{R}_A(x,y)=\frac{1}{\Pi(x)\overline{\Pi(y)}}\sum_{j=1}^{r} y_{A,j}(x)\,\overline{y_{A,j}(y)}\; }.
  \]
  \item 杠杆缺口为向量平方
  \[
  \boxed{\;\ell_A(x)=\mathcal{R}_A(x,x)=\frac{1}{|\Pi(x)|^2}\sum_{j=1}^{r}|y_{A,j}(x)|^2\; }.
  \]
  \item 对每个 \(a\in A\) 与每个 \(j\)，有 \(y_{A,j}(a)=0\)。因此存在次数 \(\le r-1\) 的多项式 \(q_{A,j}\) 使
  \[
  y_{A,j}(z)=y_A(z)\,q_{A,j}(z),
  \qquad
  y_A(z):=\prod_{a\in A}(z-a),
  \]
  并且
  \[
  \ell_A(x)=\frac{|y_A(x)|^2}{|\Pi(x)|^2}\sum_{j=1}^{r}|q_{A,j}(x)|^2.
  \]
  \item 可取 \(\{y_{A,j}\}\) 满足残差内积的正交归一性
  \[
  \boxed{\;
  \sum_{k=1}^{\kappa}\frac{y_{A,i}(p_k)\overline{y_{A,j}(p_k)}}{w_k|\Pi'(p_k)|^2}=\delta_{ij}\; }.
  \]
  并且该族在整体酉变换意义下唯一。
\end{enumerate}
\end{theorem}
\begin{proof}
记 \(U:=\Span\{v(a):a\in A\}\subset\CC^\kappa\)，则 \(\dim U=n\)，故 \(\dim U^\perp=r\)。
取 \(U^\perp\) 的一组正交归一基 \(u^{(1)},\dots,u^{(r)}\)。
令
\[
F_j(x):=\langle v(x),u^{(j)}\rangle\qquad(1\le j\le r).
\]
由投影分解
\(
\mathrm{proj}_{U^\perp}v(x)=\sum_{j=1}^r F_j(x)\,u^{(j)}
\)
得到
\[
\mathcal{R}_A(x,y)=\langle \mathrm{proj}_{U^\perp}v(x),\mathrm{proj}_{U^\perp}v(y)\rangle=\sum_{j=1}^r F_j(x)\overline{F_j(y)},
\]
从而 \(\ell_A(x)=\sum_j|F_j(x)|^2\)。
\par
每个 \(F_j\) 是在 \(\{p_k\}\) 处具有简单极点的有理函数，且无多项式部分，因此可唯一写为
\(
F_j(z)=y_{A,j}(z)/\Pi(z)
\)
其中 \(y_{A,j}\) 为次数 \(\le \kappa-1\) 的多项式。
代回得到第一、二条。
\par
由 \(u^{(j)}\in U^\perp\) 得对每个 \(a\in A\) 有 \(F_j(a)=\langle v(a),u^{(j)}\rangle=0\)，从而 \(y_{A,j}(a)=0\)。
故 \(y_{A,j}\) 被 \(y_A\) 整除，并可写为 \(y_{A,j}=y_A q_{A,j}\) 且 \(\deg q_{A,j}\le r-1\)，得到第三条。
\par
最后，比较留数：由 \(F_j(z)=\sum_k \sqrt{w_k}\overline{u^{(j)}_k}/(z-p_k)\) 与 \(F_j=y_{A,j}/\Pi\) 得
\(
y_{A,j}(p_k)/\Pi'(p_k)=\sqrt{w_k}\,\overline{u^{(j)}_k}
\)。
因此
\[
\sum_{k=1}^{\kappa}\frac{y_{A,i}(p_k)\overline{y_{A,j}(p_k)}}{w_k|\Pi'(p_k)|^2}
=
\sum_{k=1}^{\kappa}u^{(j)}_k\overline{u^{(i)}_k}
=\delta_{ij},
\]
得到正交归一性。
若更换 \(U^\perp\) 的正交归一基，则 \(u^{(j)}\) 仅差一个 \(r\times r\) 酉矩阵，故 \((y_{A,j})\) 仅差整体酉变换。证毕。
\end{proof}

\begin{proposition}[扩展 Gram 行列式的余核 Plücker 因子分解]\label{prop:xi-basepoint-scan-gram-plucker-residual-factorization}
沿用定理 \ref{thm:xi-basepoint-scan-residual-kernel-polynomial-frame} 的设定，并取互异实点
\(
X=\{x_1,\dots,x_r\}\subset\RR\setminus A
\)。
则扩展 Gram 行列式满足
\[
\boxed{\;
\frac{\det G_{A\cup X}}{\det G_A}
=\det\bigl(\mathcal{R}_A(x_i,x_j)\bigr)_{1\le i,j\le r}\; }.
\]
进一步地，若 \(y_{A,j}=y_A q_{A,j}\) 如定理 \ref{thm:xi-basepoint-scan-residual-kernel-polynomial-frame}，
则有完全显式因子化
\[
\boxed{\;
\det\bigl(\mathcal{R}_A(x_i,x_j)\bigr)_{i,j}
=
\frac{\prod_{i=1}^{r}|y_A(x_i)|^2}{\prod_{i=1}^{r}|\Pi(x_i)|^2}\cdot
\left|\det\bigl(q_{A,j}(x_i)\bigr)_{1\le i,j\le r}\right|^2\; }.
\]
右端与正交归一基的选择无关。
\end{proposition}
\begin{proof}
将 \(G_{A\cup X}\) 写成分块矩阵
\[
G_{A\cup X}=
\begin{pmatrix}
G_A & K(A,X)\\
K(X,A) & K(X,X)
\end{pmatrix}.
\]
由 Schur 补行列式公式得
\[
\det G_{A\cup X}=\det G_A\cdot \det\!\Bigl(K(X,X)-K(X,A)G_A^{-1}K(A,X)\Bigr).
\]
括号内矩阵的 \((i,j)\) 元即为 \(\mathcal{R}_A(x_i,x_j)\)，得第一式。
\par
由定理 \ref{thm:xi-basepoint-scan-residual-kernel-polynomial-frame}，矩阵 \(\bigl(\mathcal{R}_A(x_i,x_j)\bigr)_{i,j}\) 可写为
\(
D\,(Q Q^\ast)\,D^\ast
\)，其中
\[
D:=\mathrm{diag}\Bigl(\frac{y_A(x_1)}{\Pi(x_1)},\dots,\frac{y_A(x_r)}{\Pi(x_r)}\Bigr),
\qquad
Q:=(q_{A,j}(x_i))_{1\le i,j\le r}.
\]
因此
\(
\det(\mathcal{R})=|\det D|^2\det(QQ^\ast)=|\det D|^2|\det Q|^2
\)，得到第二式。
正交归一基更换对应 \(Q\mapsto Q U\)（\(U\) 酉），故 \(|\det Q|\) 不变。证毕。
\end{proof}

\begin{conjecture}[残差多项式范数的统一衰减机制]\label{conj:xi-basepoint-scan-residual-polynomial-norm-decay}
设给定一族以 \(m\) 为索引的有限秩核数据 \(\{(p_k^{(m)},w_k^{(m)})\}_{k\le \kappa_m}\)（其中 \(p_k^{(m)}\in\CC_-\)、\(w_k^{(m)}>0\)）并诱导核 \(K^{(m)}\) 与余核 \(\mathcal{R}_{A}^{(m)}\)。
则存在固定余维数 \(r\ge 1\) 与一族锚点选取规则 \(A_m\subset\RR\)（\(|A_m|=\kappa_m-r\) 且 \(G_{A_m}^{(m)}\) 可逆）使得
\[
\sup_{x\in\RR}\ \frac{\ell_{A_m}^{(m)}(x)}{K^{(m)}(x,x)}\ \longrightarrow\ 0\qquad(m\to\infty),
\]
其中 \(\ell_{A_m}^{(m)}(x)=\mathcal{R}_{A_m}^{(m)}(x,x)\)。
等价地，令 \(y_{A_m}(z)=\prod_{a\in A_m}(z-a)\)、\(\Pi_m(z)=\prod_{k\le\kappa_m}(z-p_k^{(m)})\)，并取定理 \ref{thm:xi-basepoint-scan-residual-kernel-polynomial-frame} 的残差商多项式 \(q_{A_m,j}^{(m)}\)（\(1\le j\le r\)），则上式可被视为
\[
\sup_{x\in\RR}\ \frac{|y_{A_m}(x)|^2}{|\Pi_m(x)|^2}\sum_{j=1}^{r}\bigl|q_{A_m,j}^{(m)}(x)\bigr|^2\ \longrightarrow\ 0
\]
在核化口径下的统一表述。
\end{conjecture}

\endinput


\begin{corollary}[余维一“最后一步”\texorpdfstring{$D$}{D}-最优补全：目标函数的权重无关性]\label{cor:xi-basepoint-scan-doptimal-last-step-weight-free}
固定任意 \((\kappa-1)\) 点锚集 \(A\subset\RR\)（互异）并假设 \(G_A\) 可逆。
则在所有 \(x\in\RR\setminus A\) 中最大化 \(\det G_{A\cup\{x\}}\) 的问题等价于
\[
\boxed{\;
x\in\arg\max_{t\in\RR\setminus A}\ \left|\frac{y_A(t)}{\Pi(t)}\right|\; }.
\]
因此“最后一步”最优补全点的集合完全由 \((A,\{p_k\})\) 决定，与权重 \(\{w_k\}\) 无关。
\end{corollary}
\begin{proof}
由命题 \ref{prop:xi-basepoint-scan-gram-determinant-schur-increment}，
\(
\det G_{A\cup\{x\}}=\det G_A\cdot \ell_A(x)
\)，其中 \(\det G_A\) 与 \(x\) 无关。
当 \(|A|=\kappa-1\) 时，定理 \ref{thm:xi-basepoint-scan-leverage-gap-codim1-closed-form} 给出
\(
\ell_A(x)=C(A,P,w)\cdot\bigl|\frac{y_A(x)}{\Pi(x)}\bigr|^2
\)
且 \(C(A,P,w)\) 与 \(x\) 无关。
故最大化 \(\det G_{A\cup\{x\}}\) 等价于最大化 \(\left|\frac{y_A(x)}{\Pi(x)}\right|\)。证毕。
\end{proof}

\begin{theorem}[余维一精确贪婪补全的权重独立性]\label{thm:xi-basepoint-scan-codim1-exact-greedy-weight-independence}
\leanverified{paper\_xi\_basepoint\_scan\_codim1\_exact\_greedy\_weight\_independence}
设 \(|A|=\kappa-1\) 且 \(G_A\) 可逆。
则所有使
\[
x\longmapsto \det G_{A\cup\{x\}}
\]
达到全局最大值的实点集合，与权重 \(\{w_k\}_{k=1}^{\kappa}\) 无关；它仅由锚集 \(A\) 与极点集合 \(\{p_k\}_{k=1}^{\kappa}\) 决定。
更具体地，任一极大点 \(x_\ast\in\RR\setminus A\) 都满足权重无关的平衡方程
\[
\boxed{\;
\sum_{a\in A}\frac{1}{x_\ast-a}
=
\Re\sum_{k=1}^{\kappa}\frac{1}{x_\ast-p_k}.\;}
\]
\end{theorem}
\begin{proof}
由推论 \ref{cor:xi-basepoint-scan-doptimal-last-step-weight-free}，
最大化 \(\det G_{A\cup\{x\}}\) 等价于最大化
\(
\bigl|\frac{y_A(x)}{\Pi(x)}\bigr|
\)，
故极大点集合不依赖于权重 \(\{w_k\}\)。

另一方面，由推论 \ref{cor:xi-basepoint-scan-leverage-gap-codim1-critical-points}，
任意内点临界值都满足
\[
\sum_{a\in A}\frac{1}{x-a}
=
\Re\sum_{k=1}^{\kappa}\frac{1}{x-p_k}.
\]
而在余维一情形 \(m=1\) 时，定理 \ref{thm:xi-basepoint-scan-leverage-gap-universal-xminus2-tail} 给出
\(
\ell_A(x)=c_Ax^{-2}+O(x^{-3})
\)，
故 \(\ell_A(x)\to 0\) 当 \(|x|\to\infty\)。
又由定理 \ref{thm:xi-basepoint-scan-leverage-gap-codim1-closed-form}，对 \(a\in A\) 有 \(\ell_A(a)=0\)。
因此全局极大值只能在 \(\RR\setminus A\) 的某个有界连通分支内作为内点取得，从而必满足上述临界方程。证毕。
\end{proof}

\endinput

\endinput
